\chapter{Panel data}
\label{cap:painel}

\section{The unobserved effects model}
We have $N$ units (people, firms, municipalities) observed over $T$ periods:
\[
y_{it}=x_{it}\beta+c_i+u_{it}.
\]
$c_i$ is the unobserved effect, constant over time (for example, ability or firm culture). The central question is: \textbf{is $c_i$ correlated with $x_{it}$?}

\section{The main estimators}
\begin{center}
\small
\begin{tabular}{p{3.4cm}p{5.3cm}p{5.2cm}}
\toprule
\textbf{Estimator} & \textbf{What it does} & \textbf{Consistent if} \\
\midrule
Pooled OLS (POLS) & Ignores $c_i$ and runs OLS on everything & $\E(x_{it}'(c_i+u_{it}))=0$, that is, $c_i$ uncorrelated with $x$ \\
Fixed effects (within) & Subtracts each unit's mean, eliminating $c_i$ & Strict exogeneity: $u_{it}$ uncorrelated with $x_{is}$ in \emph{all} periods \\
First differences & Regresses $\Delta y_{it}$ on $\Delta x_{it}$ & Same; with $T=2$ it coincides with fixed effects \\
Random effects (GLS) & Subtracts a fraction $\theta$ of the unit mean & $c_i$ uncorrelated with $x_{it}$ (strong assumption) \\
\bottomrule
\end{tabular}
\end{center}

\begin{intuicao}
Fixed effects uses only the variation \emph{within} each unit over time. Everything that is fixed within the unit (sex, place of birth, culture) is automatically neutralized, observed or not. The price: you lose the ability to estimate the effect of any variable that does not change over time, and you lose variation if $x$ changes little.
\end{intuicao}

\begin{simulacao}[Checked by simulation: correlated effect]
$N=500$, $T=6$, $y_{it}=1+1.0\,x_{it}+c_i+u_{it}$, with $x_{it}=0.7c_i+\text{noise}$ (that is, $c_i$ and $x$ are correlated). The true effect is 1.0.
\begin{center}
\begin{tabular}{lcccc}
\toprule
& Pooled OLS & Random effects & Fixed effects & First diff.\\
\midrule
Estimate & 1.47 & 1.27 & 0.98 & 1.00 \\
\bottomrule
\end{tabular}
\end{center}
Pooled OLS and random effects are biased because they assume $c_i$ is not correlated with $x$. Fixed effects and first differences recover the true value.
\end{simulacao}

\section{Fixed or random effects? The Hausman test and the Mundlak alternative}
The Hausman test compares fixed and random effects. A more flexible version, which in practice is more used, is the \textbf{Mundlak approach}: include the time average $\bar x_i$ as an additional regressor and test whether the coefficient of $\bar x_i$ is zero. If it is zero, random effects is reasonable. The coefficient of $\bar x_i$ measures the difference between the ``between-unit'' and ``within-unit'' effects.

\begin{simulacao}[Checked by simulation]
With $c_i$ correlated with $x$, the coefficient of $\bar x_i$ was $1.06$, with a p-value practically zero: random effects is rejected. With $c_i$ independent of $x$, the coefficient was $0.16$ (p-value $0.26$): not rejected. The value $\approx1.06$ matches the theory: the difference between the ``between'' effect ($\approx2.07$) and the ``within'' effect (1).
\end{simulacao}

\begin{atencao}[Pitfalls in panels]
\begin{itemize}[leftmargin=1.2em]
\item \textbf{The classic Hausman test under robust/clustered errors} no longer has the expected distribution. Prefer the regression version (Mundlak) with clustered errors.
\item \textbf{Always cluster the standard errors by unit.} $c_i$ generates serial correlation within the unit, even when $u_{it}$ is independent. The usual pooled OLS standard error is typically too small.
\item \textbf{Strict exogeneity is strong.} Lagged dependent variables as regressors, or \emph{feedback} ($y_{it}$ affects $x_{i,t+1}$), violate it. In those cases, dynamic panel methods (GMM) are needed, which go beyond this chapter.
\item \textbf{Fixed effects do not solve endogeneity that varies over time.}
\end{itemize}
\end{atencao}

\begin{exercicio}
With $T=2$, show that the fixed effects estimator equals the first-differences estimator.
\end{exercicio}
\sol{With $T=2$, the deviation from the mean is $\ddot y_{i1}=(y_{i1}-y_{i2})/2=-\Delta y_i/2$ and $\ddot y_{i2}=\Delta y_i/2$ (and the same holds for $x$). Then $\sum_t\ddot x_{it}\ddot y_{it}=\tfrac12\Delta x_i\Delta y_i$ and $\sum_t\ddot x_{it}^2=\tfrac12(\Delta x_i)^2$. The ratio of the sums over $i$ is $\sum\Delta x_i\Delta y_i/\sum(\Delta x_i)^2$, which is exactly OLS in first differences.}

\begin{exercicio}
Why can the effect of ``years of schooling'' not be estimated by fixed effects in a sample of adults who have already finished school?
\end{exercicio}
\sol{Because schooling does not vary within each person over the panel. After subtracting the person's mean, the regressor becomes zero in all periods, and the coefficient is not identified.}
